\section{Conclusion}
\label{sec:conclusion}

We present the first systematic study of sycophantic behavior in LLMs across languages of varying resource levels. Testing \gptlarge and \gptsmall on seven languages with two complementary experiments---multilingual nonsense engagement and factual QA capitulation---we find that sycophancy scales inversely with language resource level. The smaller model shows a statistically significant negative correlation between resource level and nonsense engagement ($\rho = -0.786$, $p = 0.036$), with engagement rates rising from 86.5\% in English to 99.0\% in Yoruba. Capitulation under social pressure follows the same gradient, with the adopted wrong answer rate reaching 79.2\% for Yoruba versus 9.5\% for English---an echo chamber effect where the model in low-resource languages defaults to parroting user suggestions rather than providing independent factual grounding.

Frontier-scale models (\gptlarge) are more robust across languages but still show elevated nonsense engagement in low-resource settings, indicating that scale alone does not close the sycophancy gap. These findings extend the multilingual alignment literature beyond safety violations to the subtler but pervasive domain of everyday reliability, with direct implications for the billions of users who interact with LLMs in non-English languages.

Future work should expand the language set (15--20 languages) and item counts (200+ per experiment) for greater statistical power, test additional model families including open-weight multilingual models, and pursue mechanistic analysis of how sycophancy and factual representations interact in activation space across languages. Professional human translations for low-resource languages and user studies in affected communities would further validate and extend these findings.
